\chapter{A \textit{ConfigManager}}

Ahogy azt a~\ref{intro-config-manager} fejezetben említettem, a \textit{ConfigManager} szoftver komponens felel a a
konfigurációs értékek beolvasásáért. Ebben a fejezetben kitérek, arra hogy hogyan értelmezi a keretrendszer a
konfigurációt, a kód oldalán és arra is hogy milyen sémájú konfigurációt vár.

\section{A konfigurációs fájl}

    A konfigurációs formátum \textit{.INI} fájl formátumot használ. Jogosan merül fel az olvasóban a kérdés, hogy miért
    pont erre a napjainkban már kevésbé divatos formátumra esett a választás.

    Az \textit{.INI} fájloknak több előnye is van, gyorsan beolvasható, a Python standard függvénykönyvtár is ad ehhez
    egy implementációt~\cite{pythonconfigparser}. Továbbá, emberek számára is olvasható, de a legfontosabb tulajdonsága
    ami miatt rá esett a választásom, hogy szekciókra lehet bontani az adatokat. Ez a szekciókra bontás lehetővé teszi
    hogy felállítsunk egy hierarchiát a konfigurációk között, így elkerülve a duplikált konfigurációkat.

    A hierarchia gyökere a \textit{DEFAULT} szekció, minden más szekció ennek az értékeit tudja felülírni vagy
    kibővíteni. A konfigurációk előtti \textit{+} jelzi hogy a szekció kibővíti az adott konfigurációs értéket, nem
    pedig felülírja azt.

    A~\ref{fig:section-hierarchy} ábra mutatja, hogyan lehet egy meglévő szekcióból konfigurációt örökölni, illetve
    felülírni. Ha a program futás során a \textit{DEFAULT} szekciót választjuk akkor a konfigurációs elemeink értékei
    rendre \verb|CFG1 = [1, 2]| és \verb|CFG2 = [3, 4]|. Ha a \textit{MY\_SECTION} szekciót választjuk akkor pedig,
    \verb|CFG1 = [1, 2, 3]| és \verb|CFG2 = [5]|.

    \begin{figure}
        \centering
        \begin{minted}{INI}
[DEFAULT]
CFG1 =
    1
    2
CFG2 =
    3
    4
[MY_SECTION]
+CFG1 =
    3
CFG2 =
    5
        \end{minted}
        \caption{\label{fig:section-hierarchy} A szekció hierarchia konstrukció}
    \end{figure}

    \pagebreak
    \subsection{A \textit{BuildOption} konfigurációk}

        A konfigurációs fájlban pár beállítástól eltekintve, a konfigurációs elemek egy-egy fordítóprogram kapcsolóhoz
        tartoznak. Ezek a \textit{BuildOption} beállítások. Ezek a beállítások egy leképezést alkotnak a
        konfigurációból a Python kódba.

        A támogatásukhoz fel kell venni egy-egy osztályt ami kezeli a konfigurációból beolvasott érték, fordítási
        kapcsolóra való konvertálását. Ezzel gyakorlailag függetleníteni lehet a konfigurációs értékek szintaxisát a
        fordítóprogram által várt parancssori paraméterek szintaxisától. Pár példa a teljeség igénye nélkül
        \textit{BuildOption} kapcsolóra: \verb|ARCH, OPT_LEVEL, EXCEPTIONS, BUILTIN, FLOAT_ABI|.

    \subsection{A \textit{Suites} konfiguráció}

        Az előbbiekben láttuk, hogy hogyan lehet kapcsolókat definiálni rugalmas módon a fordítóprogram számára, de
        hogyan fogja tudni hogy milyen teszteket futtasson?

        Erre ad megoldást a \textit{Suites} konfigurációs mező ahol fel kell sorolni a futtatni kívánt teszt
        csomagokat. A teszt csomagokhoz tartozó fájlokat egy \textit{SET} forrásfájlban kell felsorolni, minden sorban
        egy elérési útvonallal. Ezzel a mechanizmussal bármilyen teszt esetekből álló teszt csomag definiálható.

        A legnagyobb előnye ennek a formátumnak, hogy a keretrendszer követni tudja, hogy melyik teszt esetet futtatta
        már. Így ha a teszt csomagok osztoznak egy-egy teszt eseten akkor a keretrendszer egyszerűen
        \textbf{behelyettesíti} a meglévő eredményt.

        Ezen mechanizmus adta az ötletet a projekt elnevezésére.

    \subsection{A \textit{Matrix} és a \textit{Constraints}}

        Most, hogy már ki tudjuk választani a fordítóprogram számára megfelelő opciókat, és definiálni tudjuk a
        futtatandó teszt csomagokat a kettő beállítás között kapcsolatot kell teremteni. Erre a kapcsolat teremtésre
        szolgálnak a \textit{Matrix} és a \textit{Constraints} beállítások.

        A \textit{Matrix} beállításban fel kell sorolni az összes olyan \textit{BuildOption} konfigurációt amit
        futtatni szeretnénk.

        A \textit{Constraints} beállításban adhatunk meg megszorításokat, hogy bizonyos teszt csomagokra csak bizonyos,
        kapcsolókat engedélyezzünk. Így a \textit{Matrix} által generált potenciálisan exponenciális kombinációkat
        levághatóak, ha nem szükségesek. A szintaxist a következő \textit{EBNF}~\ref{fig:ebnf} nyelvtan írja le.

        \begin{figure}[h]
            \hspace*{-1.5in}
            \centering
            \begin{grammar}
                <constraint> ::= <suite> <build_option_constraint>+ %chktex 26

                <build_option_constraint> ::= "&" <id> %chktex 18 chktex 26
                \alt "&" <build_option_with_value> %chktex 1 chktex 18

                <suite> ::= "*" %chktex 18 chktex 26
                \alt <id> %chktex 1 chktex 18

                <id> ::= [a-zA-Z0-9_]+ %chktex 8 chktex 18 chktex 26
            \end{grammar}
            \caption{\label{fig:ebnf} A \textit{Constraints} kapcsoló nyelvtana}
        \end{figure}
